% =====================================================================
%  Section 4.2 — Emergent distorted spectrum
%
%  Figures (generated by spectrum_diagnostic.py):
%     fig_spec_rdt.pdf      RUN A2  : single-case, eddies off + dilatation on
%     fig_spec_cascade.pdf  RUN B   : single-case, eddies on + dilatation on
%     fig_spec_compare.pdf  RUN C   : comparison over the kvisc0 sweep
%  Requires graphicx, amsmath. Cross-refs: sec:level1a, eq:plane-strain,
%  eq:lrr, sec:val-spectrum (this section's label), sec:gate (Gate 2 / Sec 4.3).
% =====================================================================

\subsection{Emergent distorted spectrum}\label{sec:val-spectrum}

The verification of \S\,\ref{sec:level1a} establishes that the solver
reproduces the single-point moment evolution of rapid-distortion theory. That
is a necessary check, but it is not the contribution: the single-point
statistics are, by construction, those of the closure, and a closure alone
could supply them. What ODT adds, and what Amiet's theory ultimately requires,
is the \emph{spectrum} of the distorted upwash---how the incoming turbulent
energy is redistributed across wavenumber as the gust is compressed toward the
leading edge. This section examines that spectrum, built up in three steps:
the geometric compression alone (a linear-RDT baseline), the full model with
the eddy events active, and the dependence on Reynolds number.

\subsubsection{Compression without cascade: the linear-RDT
baseline}\label{sec:spec-rdt}

The mean strain compresses the line along its axis. With the events
suppressed and zero molecular viscosity, this compression is the only
process acting on the resolved field, and its effect is purely kinematic:
the line length contracts as $\mathcal{L}(e)=\mathcal{L}_0\,e^{A_{22}e}$, so
every resolved wavenumber is rescaled as $k(e)=k(0)\,e^{-A_{22}e}$ with no
transfer of energy between scales. This is precisely the linear
rapid-distortion result for the transverse wavenumber, and it provides an
exact, parameter-free target for the implementation.

Figure~\ref{fig:spec-rdt} shows the outcome. As the strain increases the
component spectra translate rigidly toward higher wavenumber---amplitude
growing anisotropically, with the upwash component $E_2$ amplified---while
the spectral shape is preserved: the peak does not broaden and no inertial
range develops. The spectral centroid migrates as $e^{-A_{22}e}$ to within
the discretisation error, reaching $7.0$ at $e=3.9$ against the exact value
$e^{-A_{22}\cdot3.9}=7.03$, and the three components migrate identically,
confirming that the shift is geometric rather than dynamical. The
single-point moments are unchanged from the no-dilatation case (the
fractions remain on the LRR trajectory of \S\,\ref{sec:level1a}), since a
uniform rescaling of position leaves the length-weighted velocity statistics
invariant. The compression therefore acts entirely on the spectrum and not
on the moments, and reproduces linear RDT exactly. It is the baseline against
which the cascade is measured.

\begin{figure}
  \centering
  \includegraphics[width=0.92\textwidth]{fig_spec_rdt}
  \caption{Compression without cascade (eddy events suppressed, inviscid).
  Left: component spectra $E_i(k)$ at increasing total strain
  (lighter to darker); the spectra translate to higher $k$ without changing
  shape. Right: spectral-centroid migration $\bar k(e)/\bar k(0)$ for each
  component (symbols) against the linear-RDT prediction $e^{-A_{22}e}$
  (dashed). The centroid follows the prediction exactly and the three
  components coincide, confirming purely geometric compression.}
  \label{fig:spec-rdt}
\end{figure}

\subsubsection{The emergent cascade}\label{sec:spec-cascade}

Activating the eddy events introduces the process the linear baseline lacks:
a turbulent cascade that transfers energy across scales as the gust is
strained. Because the events also realise the slow (return-to-isotropy)
redistribution, the single-point fractions now depart from the LRR
trajectory---expected, and not a defect of the implementation. The
quantity of interest is the spectrum.

Figure~\ref{fig:spec-cascade} contrasts the full model with the compression
baseline of Figure~\ref{fig:spec-rdt}. Where the baseline spectrum
translated rigidly, the full-model spectrum \emph{broadens}: energy is
carried to wavenumbers beyond those reached by compression alone, and the
spectral centroid rises above the geometric line. This excess over the
linear-RDT baseline is the cascade---the transfer that rapid-distortion
theory, which only rescales the prescribed spectrum, cannot represent. It is
the mechanism by which the strain-coupled model produces a distorted upwash
spectrum that differs from the frozen, linearly-strained spectrum Amiet's
theory assumes.

\begin{figure}
  \centering
  \includegraphics[width=0.92\textwidth]{fig_spec_cascade}
  \caption{The full model (eddy events on, dilatation on). The distorted
  spectrum broadens beyond the rigid translation of Figure~\ref{fig:spec-rdt},
  and the centroid rises above the linear-RDT compression line: energy is
  transferred to higher wavenumber by the cascade, beyond what geometric
  compression alone supplies.}
  \label{fig:spec-cascade}
\end{figure}

\subsubsection{Dependence on Reynolds number}\label{sec:spec-Re}

In these non-dimensional variables the molecular viscosity is the inverse
Reynolds number, $\nu=\texttt{kvisc0}\approx Re_{\mathcal{L}}^{-1}$, and it
sets the depth of the cascade---the separation between the energy-containing
peak and the dissipation scale. Figure~\ref{fig:spec-compare} sweeps it over
two decades. As $\nu$ decreases the spectrum extends progressively further
toward high wavenumber and the centroid climbs: the cascade deepens
monotonically with Reynolds number, as it must physically. This trend is the
central qualitative result---the model produces a Reynolds-number-dependent
distortion of the spectrum, a behaviour outside the reach of linear theory.

Two cautions temper the quantitative reading, and we state them plainly. At
the lowest viscosities the spectrum does not roll off cleanly but flattens
into a shelf at the highest resolved wavenumbers: the dissipation scale has
fallen below the grid, so energy piles up at the cutoff rather than
dissipating. The centroid of these cases is therefore inflated by an
under-resolved, partly numerical contribution, and they should not be read
as physical cascade depths without a resolution check. Conversely, the
adequately resolved cases do not yet exhibit an extended $k^{-5/3}$ inertial
range over this strain interval. The cascade is real and deepens with
Reynolds number, but its quantitative form at high Reynolds number is
resolution-limited here; resolving it is a matter of grid density rather than
of the model.

\begin{figure}
  \centering
  \includegraphics[width=0.92\textwidth]{fig_spec_compare}
  \caption{Dependence on Reynolds number ($\nu=\texttt{kvisc0}$, eddy events
  and dilatation on). Left: final-strain spectrum $E(k)=\sum_i E_i(k)$ for
  decreasing $\nu$; the cascade extends further toward high $k$ as $\nu$
  falls. The dotted line marks a $k^{-5/3}$ slope for reference. At the
  smallest $\nu$ the spectrum flattens into a grid-scale shelf at the cutoff
  (under-resolved). Right: centroid migration per case against the linear-RDT
  line (dashed).}
  \label{fig:spec-compare}
\end{figure}

This fixes how the operating Reynolds number must be chosen. It is a physical
property of the inflow turbulence the model represents---set by matching
$Re_{\mathcal{L}}$ of the target grid- or inflow-turbulence, hence
$\nu=Re_{\mathcal{L}}^{-1}$---and not a free parameter to be tuned until the
centroid exceeds the compression line. Once chosen physically, the value must
be checked for resolution: the dissipation rolloff must sit below the grid
cutoff with margin. The cases that most strongly exceed the linear baseline
in Figure~\ref{fig:spec-compare} are precisely those failing this check, so
the physical operating point is the largest $\nu$ (lowest $Re$) consistent
with the target inflow, resolved on the available grid. With that point
fixed, the distorted upwash spectrum produced here is the input to the
leading-edge response of \S\,\ref{sec:gate}.
